\documentclass[11pt]{article}
\usepackage[utf8]{inputenc}
\usepackage[T1]{fontenc}
\usepackage{lmodern}
\usepackage[margin=1in]{geometry}
\usepackage{amsmath,amssymb,amsthm}
\usepackage{enumitem}
\usepackage{booktabs}
\usepackage{xcolor}
\usepackage{tikz}
\usepackage[colorlinks=true,bookmarks=false,linkcolor=blue!50!black,urlcolor=blue!50!black]{hyperref}

\theoremstyle{plain}
\newtheorem{theorem}{Theorem}
\newtheorem{lemma}[theorem]{Lemma}
\newtheorem{claim}[theorem]{Claim}
\newtheorem{corollary}[theorem]{Corollary}
\theoremstyle{definition}
\newtheorem{definition}[theorem]{Definition}
\newtheorem{example}[theorem]{Example}
\theoremstyle{remark}
\newtheorem*{remark}{Remark}

\newcommand{\N}{\mathbb{N}}
\newcommand{\Z}{\mathbb{Z}}
\newcommand{\R}{\mathbb{R}}
\newcommand{\Q}{\mathbb{Q}}
\newcommand{\iv}[1]{[#1]}                 % Iverson bracket
\newcommand{\powerset}{\mathcal{P}}
\DeclareMathOperator{\ord}{ord}
\DeclareMathOperator{\gcdop}{gcd}

\setlength{\parindent}{0pt}
\setlength{\parskip}{0.6em}
\newcommand{\lesson}[2]{\href{../lessons/#1}{#2}}   % link to a lesson image

% Let long inline lists like \{1,2,5,6\} break across lines after a comma.
\mathchardef\mycomma="613B
\mathcode`,="8000
{\catcode`,=\active \gdef,{\mycomma\penalty0\hspace{0pt}}}
\emergencystretch=2em

\title{Exponential generating functions\\ and labelled structures}
\author{Notes from Carlos' lessons}
\date{}
\begin{document}
\maketitle

\textbf{Source lessons:}
\lesson{2026-06-13-exponential-generating-functions.png}{06-13 (EGFs and sums of powers)},
\lesson{2026-06-21-exponential-generating-functions-vs-algebraic-datatypes.png}{06-21 (EGFs vs.\ algebraic datatypes)},
\lesson{2026-07-07-exponential-generating-functions-multinomial-coefficients.png}{07-07 ($e^{A(x)}$ and multinomial coefficients)}.

\section{Definition: OGF vs.\ EGF (06-13)}

\begin{center}
\begin{tabular}{@{}lll@{}}
\toprule
 & map & $(1,1,1,\dots)\mapsto$\\
\midrule
ordinary (OGF) & $(a_0,a_1,\dots)\mapsto\sum a_nx^n$ & $\dfrac1{1-x}$\\[1.2ex]
exponential (EGF) & $(a_0,a_1,\dots)\mapsto\sum a_n\dfrac{x^n}{n!}$ & $e^x$\\
\bottomrule
\end{tabular}
\end{center}
(The board also recalls $e^x=\lim_{n\to\infty}(1+\frac xn)^n$ and $e=\lim(1+\frac1n)^n$.)

\section{The product rule (06-21) and its meaning}

\begin{claim}
If $A(x)=\sum a_n\frac{x^n}{n!}$ and $B(x)=\sum b_n\frac{x^n}{n!}$, then
\[
A(x)B(x)=\sum_{n\ge0}\Bigl(\sum_{k=0}^n\binom nk a_kb_{n-k}\Bigr)\frac{x^n}{n!}.
\]
\end{claim}
\begin{proof}
$[x^n]AB=\sum_{k=0}^n\frac{a_k}{k!}\cdot\frac{b_{n-k}}{(n-k)!}=\frac1{n!}\sum_k\binom nk a_kb_{n-k}$.
\end{proof}

\paragraph{Why EGFs count labelled structures.} Suppose $a_n$ is the number of structures of type $A$ on the label set $[n]$, and $b_n$ likewise for $B$. A pair (an $A$-structure, a $B$-structure) on $[n]$ is built by: (1) choosing the size $k$ of the label set of the left component; (2) choosing \emph{which} $k$ labels go left, in $\binom nk$ ways; (3) putting an $A$-structure on those labels and a $B$-structure on the remaining $n-k$, in $a_kb_{n-k}$ ways. That is exactly the coefficient above. (The board contrasts this with OGFs, where the product counts pairs with sizes adding, with no relabelling.) The sum $A+B$ counts ``an $A$ or a $B$ structure'' (the board's case classes $T_1\cup T_2$), and $A\cdot B$ counts pairs ($T_1\times T_2$) with $\mathrm{size}(t_1,t_2)=\mathrm{size}(t_1)+\mathrm{size}(t_2)$.

\subsection*{Application: the binomial theorem (06-21 board)}

Take $a_k=\alpha^k$ and $b_k=\beta^k$: $A(x)=\sum\frac{(\alpha x)^k}{k!}=e^{\alpha x}$ and $B(x)=e^{\beta x}$. Then $A(x)B(x)=e^{(\alpha+\beta)x}$ has $n$-th coefficient $(\alpha+\beta)^n$, while the product rule gives $\sum_k\binom nk\alpha^k\beta^{n-k}$. Hence
\[
(\alpha+\beta)^n=\sum_k\binom nk\alpha^k\beta^{n-k}.
\]
(Compare the combinatorial proof in the topic \emph{Counting and double counting}.)

\section{Sums of powers via EGFs (06-13)}

Let $S_k(n)=\sum_{i=1}^ni^k$ for $k\ge0$, so $S_0(n)=n$ and $S_1(n)=\frac{n(n+1)}2$.

\subsection*{Elementary (telescoping) derivation}

$(i+1)^{k+1}-i^{k+1}=\sum_{\ell=0}^k\binom{k+1}\ell i^\ell$ by the binomial theorem. Sum over $i=1,\dots,n$; the left side telescopes:
\[
(n+1)^{k+1}-1=\sum_{\ell=0}^k\binom{k+1}\ell S_\ell(n).
\]
Isolating the $\ell=k$ term ($\binom{k+1}k=k+1$) gives the recursion boxed on the board:
\[
S_k(n)=\frac{(n+1)^{k+1}-1-\sum_{\ell=0}^{k-1}\binom{k+1}\ell S_\ell(n)}{k+1}.
\]
\emph{Example ($k=2$).} $(n+1)^3-1=\binom30n+\binom31\frac{n(n+1)}2+3S_2(n)$, i.e.\ $n^3+3n^2+3n=n+\frac32n(n+1)+3S_2$, so
\[
S_2(n)=\tfrac13n^3+\tfrac12n^2+\tfrac16n.
\]

\subsection*{EGF derivation (the same identity, packaged)}

Put $F_n(x)=\sum_{k\ge0}S_k(n)\frac{x^k}{k!}$. Swapping the sums, $F_n(x)=\sum_{i=1}^n\sum_k\frac{(ix)^k}{k!}=\sum_{i=1}^ne^{ix}=\dfrac{e^{(n+1)x}-e^x}{e^x-1}$ (a geometric series in $e^x$). The board writes the version summing from $m=0$, $\frac{e^{(n+1)x}-1}{e^x-1}$, which differs only by the $i=0$ term, i.e.\ it counts $0^0=1$ in $S_0$. Multiply by $e^x-1=\sum_{j\ge1}\frac{x^j}{j!}$:
\[
F_n(x)\,(e^x-1)=e^{(n+1)x}-e^x.
\]
Take $[x^{k+1}]\cdot(k+1)!$. The left side is the EGF product of $(S_\ell(n))_\ell$ with $(0,1,1,1,\dots)$, i.e.\ $\sum_{\ell=0}^{k}\binom{k+1}\ell S_\ell(n)$, and the right side is $(n+1)^{k+1}-1$. This is the identity above. (It is the board's chain, rewritten with the EGF $e^x-1$ of $(0,1,1,1,\dots)$.)

\section{Algebraic datatypes as generating-function equations (06-21)}

A recursive datatype defines a combinatorial class: sum types become $+$, product types become $\times$, and the size of a value is the sum of the sizes of its parts. Translate directly to an OGF equation (sizes only, no relabelling needed).

\begin{itemize}
  \item \textbf{FibList.} \texttt{FibList = [] | 1 (FibList) | 2 (FibList)}, where the value $1$ or $2$ contributes size $1$ or $2$ (the board: ``finite list consisting of 1s and 2s, size $=\ell_1+\dots+\ell_k$''). Then
  \[
  F=1+xF+x^2F\ \Rightarrow\ F=\frac1{1-x-x^2},
  \]
  so the number of lists of total size $n$ is $[x^n]F=F_{n+1}$ (Fibonacci), equivalently the number of ways to climb $n$ steps taking $1$ or $2$ at a time (the staircase pictures on the board, and the definition of $F_n$ as the number of sequences $s_i\in\{1,2\}$ with $s_1+\dots+s_k=n$). Compare the topic \emph{Ordinary generating functions}.
  \item \textbf{BinaryTree.} \texttt{BinaryTree = Node Int | Tree BinaryTree Int BinaryTree}, with each \texttt{Int} of size $1$:
  \[
  B(x)=x+x\,B(x)^2.
  \]
  Write $B=x\,C(x^2)$ where $C=1+yC^2$ is the Catalan series ($y=x^2$). Indeed $x+x\cdot x^2C^2=x(1+x^2C^2)=xC$. So
  \[
  B(x)=\frac{1-\sqrt{1-4x^2}}{2x}=\sum_kC_kx^{2k+1}:
  \]
  the number of such trees with $2k+1$ nodes is the Catalan number $C_k$. (The board's tree pictures show the size-$5$ case.)
\end{itemize}

\section{$e^{A(x)}$: sets of structures (07-07)}

Given an EGF $A(x)=\sum_{k\ge1}a_k\frac{x^k}{k!}$ with $a_0=0$, the board studies $e^{A(x)}=\sum_{n\ge0}\frac{A(x)^n}{n!}$ and extracts its coefficients with two rules:
\begin{enumerate}[nosep]
  \item $[x^i]A(x)^n=\sum_{i_1+\dots+i_n=i}[x^{i_1}]A\cdots[x^{i_n}]A$ (expand the product),
  \item $[x^i]C(x)D(x)=\sum_{s_1+s_2=i}[x^{s_1}]C\,[x^{s_2}]D$.
\end{enumerate}
Therefore
\[
[x^i]\,e^{A(x)}=\sum_{n\ge0}\frac1{n!}\sum_{i_1+\dots+i_n=i}\frac{a_{i_1}\cdots a_{i_n}}{i_1!\cdots i_n!}.
\]

\paragraph{Combinatorial meaning.} $A(x)^n$ counts ordered $n$-tuples of $A$-structures on a labelled set partitioned into $n$ non-empty blocks. The blocks have \emph{distinct} label sets, so reordering produces $n!$ different ordered tuples of the same unordered collection; dividing by $n!$ counts \textbf{sets} of $n$ structures. So $e^{A(x)}$ is the EGF of ``a set of $A$-structures'' (with no restriction on how many).

The board's picture of the same principle:
\begin{itemize}[nosep]
  \item $A=(1,1,1,\dots)$ has EGF $e^x$ (``sets'': one structure on each label set), and $e^x-1$ is ``non-empty sets''.
  \item Two applications are in the topic \emph{Set partitions and involutions}: $e^{e^x-1}$ (set partitions) and $e^{x+x^2/2}$ (involutions).
\end{itemize}

\section{Binomial and multinomial coefficients (07-07, right side)}

\begin{itemize}[nosep]
  \item $\binom nk=\#\{k\text{-subsets of }\{1,\dots,n\}\}=\binom n{k,\,n-k}$ (``$A$ and its complement'').
  \item \textbf{Multinomial.} $\binom n{k_1,\dots,k_p}$ (with $k_1+\dots+k_p=n$) counts the ordered partitions of $\{1,\dots,n\}$ into blocks of sizes $k_1,\dots,k_p$.
\end{itemize}
\begin{proof}
Choose the first block, then the second from what remains, and so on:
\[
\binom n{k_1}\binom{n-k_1}{k_2}\binom{n-k_1-k_2}{k_3}\cdots\binom{k_p}{k_p}
=\frac{n!}{k_1!(n-k_1)!}\cdot\frac{(n-k_1)!}{k_2!(n-k_1-k_2)!}\cdots=\frac{n!}{k_1!\cdots k_p!},
\]
because the intermediate factorials cancel (the slash marks on the board) and $\binom{k_p}{k_p}=1$.
\end{proof}

\section*{Notes}
\begin{itemize}[nosep]
  \item The right half of the 06-13 board also lists $\sum_n\frac{e^{(n-1)x}-1}{e^x-1}$ as an alternative form of $F_n(x)$ (consistent up to an index shift) and the product/EGF rewriting $E(a_k)\cdot E(b_k)=\sum\frac{x^n}{n!}\sum_k\binom nka_kb_{n-k}$.
  \item The large grid sketch at the upper left of the 06-21 board (curves and a staircase between $n$ and $h$) is a staircase/lattice picture that I could not interpret with confidence; I have only transcribed the algebra beside it.
\end{itemize}

\end{document}
